\section{Limits of fixed-copy scalar invariants}\label{sec:invariants}

The operator-valued marginal mechanism retains local Weyl axes.  Fixed-copy
scalar contractions retain less: on algebraic families they are generically
constant and can distinguish parameters only on rank-jump strata.  We make
that limitation precise here; Appendix~\ref{sec:scalar-certificates}
records two explicit separators on exceptional strata.

For \(r\) ket copies and \(r\) bra copies, choose a permutation
\(\sigma_i\in S_r\) at each party and contract ket copy \(a\) with bra
copy \(\sigma_i(a)\).  If \(G\) generates a \(d\)-dimensional code, let
\(M_{\boldsymbol\sigma}(G)\) be the homogeneous matching system
\[
 (u_aG)_i=(v_{\sigma_i(a)}G)_i
 \quad(1\leq i\leq n,\ 1\leq a\leq r).
\]
Counting its solutions gives the normalized contraction
\begin{equation}\label{eq:permutation-contraction}
 I_{\boldsymbol\sigma}(\Psi_C)
   =q^{dr-\rank M_{\boldsymbol\sigma}(G)}.
\end{equation}

The next statement uses a small amount of algebraic-geometric language.  An
irreducible parameter variety is a family not decomposed into smaller closed
families, and a dense open subset is what remains after excluding the common
zero locus of finitely many nonzero polynomial conditions.  Thus
``generically constant'' means constant away from such rank-jump conditions.

\begin{proposition}[Generic constancy of fixed-copy invariants]\label{thm:fixed-copy-boundary}
Fix \(r\) and a field \(k\) of cardinality \(q\), which is also the local
Hilbert-space dimension under the code--state dictionary.  Fix an
irreducible reduced parameter variety \(X\) over \(k\), and a regular
generator-matrix chart \(G(x)\) of constant code dimension.  Every
degree-\((r,r)\) permutation contraction of the associated equal-phase
states has one value at every \(k\)-point of some dense open subscheme
\(U\subseteq X\) defined over \(k\).  The same holds for party-orbit sums and
linear combinations.  After evaluating the states as vectors in
\((\mathbb C^q)^{\otimes n}\), in the stable range \(q\geq r\), every
scalar LU-invariant of bidegree \((r,r)\) is generically constant in this
sense.  For a reducible family, the statement holds separately on each
irreducible component.
\end{proposition}

\begin{proof}
Formula~\eqref{eq:permutation-contraction} depends only on the rank of a
matrix whose entries are regular functions over \(k\).  A maximal nonzero
minor shows that this rank is constant on a dense open subscheme defined
over \(k\).  There are only finitely many diagrams at fixed \(r\), so
their generic-rank opens have a common dense intersection because \(X\)
is irreducible.  This proves the geometric generic-rank assertion.  Over the complex vector space
\((\mathbb C^q)^{\otimes n}\), applied pointwise to each code state, the
first fundamental theorem for unitary invariants says that local
copy-permutation diagrams span the bidegree-\((r,r)\) invariants; when the
local Hilbert-space dimension \(q\) is at least \(r\), there are no
additional dimension relations among these diagrams
\cite[Section~2]{CollinsSniady2006}.
\end{proof}

\begin{remark}[Finite-field points]
The dense open subscheme \(U\) need not contain a \(k\)-rational point.  If
\(U(k)=\varnothing\), the corresponding pointwise statement over the finite
field is vacuous; the geometric generic-rank statement remains valid.
\end{remark}

Thus no fixed copy number recovers a nonconstant pencil coordinate as a
regular function on the generic locus.  This is a statement about the
geometric generic locus, not about the distinguishability of individual
code states: polynomial LU invariants separate inequivalent states, so
any finite collection of pairwise inequivalent code states is separated at
some finite copy number.  What remains open is whether that copy number
can be bounded uniformly as \(q\) or the family grows.  A scalar witness
can still detect a special
rank-jump stratum, as the four-copy example in
Appendix~\ref{sec:scalar-certificates} does.  In contrast, the four-party
reduced-operator family \eqref{eq:css-reduced-state} retains the Weyl axes,
and the imported rigidity theorem forces Cliffordness.  This is the precise
distinction between generic constancy of scalar invariants and the
operator-valued marginal information used by the rigidity theorem.
